\chapter{Conclusions}\label{chapter:conclusions}

This thesis addressed the recurring infrastructure work between ROS~2
communication and property-specific checking. It reviewed the domain, developed
a feature model and reference architecture, implemented the selected path, and
evaluated five increasingly distributed experiments.

\section{Answers to the Research Questions}

\textbf{RQ1: What monitoring capabilities and ROS~2 observation mechanisms
exist, and which application-level observation gaps remain in the reviewed
approaches?}

ROS~2 behaviour can be observed through native typed interfaces,
instrumentation, network or middleware probes, and system tracing. Topic
observation is common in the reviewed ROS-focused work, ROSMonitoring 2.0 adds
services with partial ROS~2 support, and the reviewed approaches do not combine
ROS~2 topic messages, service events, and action feedback/status in one
configurable application-level infrastructure.

\textbf{RQ2: What are the commonalities and variabilities of monitoring
solutions for ROS~2 applications?}

Solutions share collection, representation, processing, routing, checking, and
result-use tasks. They vary in observable elements, observation level,
monitoring mode, online/offline operation, placement, property formalism,
transport, output, and whether artifacts are configured, generated, or
extended through plug-ins. The feature model makes these choices explicit.

\textbf{RQ3: How can a reference architecture incorporate the selected
monitoring features while separating reusable infrastructure from
property-specific checking?}

A collector converts a ROS~2 observation to a common record. Reusable
transformers and dispatchers reduce and route it. A converter plug-in adapts the
record to a checker-specific DSL record, and a verdict-service plug-in evaluates
the property. Common verdict and evidence formats retain attribution. The same
roles can run together or be connected through a configured transport.

\textbf{RQ4: To what extent do the proposed reference architecture and its
implementation support the monitoring needs of representative ROS~2
applications?}

Within the tested scope, the support is demonstrated. The experiments covered
topics, service request/response events, action feedback/status, local and
MQTT-transported checking, multi-source properties, two namespaces, and a
three-computer deployment. Reported verdicts matched controlled expectations,
replay, mission logs, or Gazebo evidence. Runtime YAML was generated from
deployment specifications, and property logic was reused unchanged when E4 was
moved to the E5 placement.

\section{Main Contribution}

The main contribution is a clear architectural and implementation boundary
between heterogeneous ROS~2 observations and checker-specific property logic.
The common record preserves observation context, while configuration and
plug-in interfaces allow sources, processing, checking, outputs, and placement
to vary. Configuration generation reduces low-level routine boilerplate without
hiding application and property decisions.

The result is intentionally a foundation rather than a universal verifier.
Active feedback, complete action traces, lower-level observation, globally
ordered distributed checking, and production-scale fault tolerance remain
outside the prototype. The evaluation establishes feasibility and functional
correctness for the reported configurations.

\section{Availability}

The implementation, tests, documentation, and case-study material are publicly
available; Appendix~\ref{app:repository} gives the repository location and
describes its contents.
